\subsection{正次元必要片の対角切片と五つの曲線表示}
\label{subsec:cubic-four-simple-positive-rank}
前節で残った75片にも対角torusを適用する。
次の曲線は必要被覆であり、AS4性や正常三次の存在をまだ保証しない。
\begin{proposition}\label{prop:cubic-four-simple-five-curves}
四単純H3包絡は、証明書の五つの一変数表示で必要被覆される。
そのスペクトルラベルは
$(5;1,2,3,4)$、$(6;0,1,2,3)$、$(7;1,2,5,6)$、$(9;1,3,6,8)$
の四つ、計9根集合に限られる。
\end{proposition}
\begin{proof}
各片の係数重み行列 $C$ の独立列を取り、それらに対応するorbit係数を1と置く。
独立な指標が定める代数閉体上のtorus写像は全射なので、全係数軌道がこの切片に交わる。
このgauge式を元の整数行列 $A$ に追加し、再びSmith分解の全有限成分を列挙する。
自由座標と全位相を保持して重複表示を除くと16表示となる。
各表示のorbit係数は原始根 $z=\zeta_L$ と非零変数 $t_j$ のLaurent単項式であり、
全75片からの対応と指数は証明書に保存する。

この係数を元の四次tensorへ戻し、各weight blockに必要rankを超える全minorを課す。
非零minorの単項式因子はtorus上の単元なので除く。
三次元表示では、正規化した二つのminorの差が $-4z^3$ となり、空である。
それ以外では全minorが次表の単一式の倍数となり、逆にその式自体が正規化minorに現れる。
従って表はrank上界の必要条件だけの一般点近似ではなく、その全零点集合を与える。
\[
\begin{array}{c|r|l}
\text{表示番号}&L&\text{rank上界の式}\\\hline
0&10&t_1t_2^2+zt_1-2t_2^2\\
1,2&20&t_1+2z^4\\
3&10&t_1^2t_2+2z^2t_1-z^2t_2\\
4&20&t_1+2z^7\\
5&20&t_1-2z^7\\
6&10&\text{空}\\
7&6&\text{全ての }t_1\ne0\\
8,9&12&t_1^2+2z^2\\
10&14&t_1t_2^2-z^2t_1-2z^3t_2\\
11&28&t_1-2z^9\\
12&28&t_1+2z^9\\
13&18&t_1^2t_2-2t_1+(z^4-z)t_2\\
14&36&t_1+2z^{11}\\
15&36&t_1-2z^{11}
\end{array}
\]
番号0,3,10,13はそれぞれ一方の変数に一次であり、その係数が零なら定数項が非零となる。
従って非零変数の範囲で各々一つの有理曲線表示となる。
番号7を加えて五つの一変数表示を得る。

他の表示は計12個のパラメータ点へ縮む。
一次式は唯一の根を取り、番号8,9の二次式は両根を取る。
後者は $w=\zeta_{24}$、$z=w^2$ として
$t_1=\pm(w^{11}+w^5)$ と表示できる。
各点の元middle像はrank6、次数3,4の商次元は20,35であるが、
正常三次行列は全て列full rankとなる。
具体的には番号1,2は $21\times4$、4,5は $23\times4$、
8,9の全四点は $26\times4$、11,12は $14\times2$、14,15は $12\times2$ である。
従って全12点で $W_1(B)=0$ となりH3を与えない。
計算は全て標数0の数体上であり、有限標本から特殊値全体を推論していない。
\end{proof}
再現器は\\
\texttt{verify\_cubic\_four\_simple\_positive\_rank.py}、\\
\texttt{verify\_cubic\_four\_simple\_positive\_finite.py}。
対応する二つのJSON証明書に全切片、全位相、全rank式と全有限点の正常rankを保持する。
五曲線の全正常三次incidence、全CY領域、成分・群oid・相対降下はこの命題では解決していない。
